\documentclass[12pt,a4paper,landscape]{article}
\usepackage[T1]{fontenc}
\usepackage{amssymb}
\usepackage{xcolor}
\usepackage{circuitikz}
\usepackage{tikz}


\title{Block Diagram}
\author{RB}


\begin{document}
	
	\section{Synchronous FIFO}
	\vspace{2cm}
	\begin{circuitikz}
	    \ctikzset{multipoles/thickness=3}
	    \ctikzset{
	        multipoles/dipchip/width=2,
	        multipoles/dipchip/pin spacing=0.75
        }
	    \draw (0,0) 
	        node[
	            dipchip, 
	            num pins=12,
	            no topmark, 
	            external pins width=0.0,
	            %hide numbers
	        ](fifo){FIFO};
	    
	    
	    % // ---------- Inputs ---------- \\
	    \draw[<-] (fifo.pin 5) -- ++(-4,0) node[left]{\small reset\_n};
        \draw[<-] (fifo.pin 6) -- ++(-4,0) node[left]{\small clk};
        
	    % -- Writes --
	    \draw[line width=2, <-] (fifo.pin 1) to[multiwire, l_=\tiny FIFO\_DATAWIDTH] ++(-4,0) node[left]{\small wdata};
	    \draw[<-] (fifo.pin 3) -- ++(-4,0) node[left]{\small wen};
	    
	    % -- Reads --
	    \draw[<-] (fifo.pin 4) -- ++(-4,0) node[left]{\small ren};
	    
	    % -- Flags --

        % // ---------- Outputs ---------- \\
        % -- Writes --
        \draw[->, line width=2] (fifo.pin 10) to[multiwire, l=\tiny FIFO\_PTR] ++(4,0) node[right]{room\_avail};
	    % -- Reads --
	    \draw[->, line width=2] (fifo.pin 12) to[multiwire, l=\tiny FIFO\_DATAWIDTH] ++(4,0) node[right]{\small rdata};
	    \draw[->, line width=2] (fifo.pin 11) to[multiwire, l=\tiny FIFO\_PTR] ++(4,0) node[right]{data\_avail};
	    % -- Flags --
	    \draw[->] (fifo.pin 7) -- ++(4,0) node[right]{\small full};
	    \draw[->] (fifo.pin 8) -- ++(4,0) node[right]{\small empty};
        
        
        
	\end{circuitikz}
	
	\newpage
	\section{Asynchrnous FIFO}
	
\end{document}